\paragraph{Source-level interpretation of the native pair.}
For the retained round-1 source \texttt{07ba285c6b6c}, let $x_6,x_{24}$ be
the fixed preprocessed protein inputs, not raw fold changes. The implemented
protein representation is
\[
\begin{aligned}
z_t&=\operatorname{ReLU}(W x_t+b),\qquad
d=\operatorname{ReLU}(V(x_{24}-x_6)+c),\\
g&=\operatorname{sigmoid}(G[z_6,z_{24},d]+a),\\
h&=g\odot(z_6+z_{24})/2+(1-g)\odot d.
\end{aligned}
\]
The shared per-time and difference encoders each produce 32 features;
the gate maps 96 to 32. The unchanged drug branch produces another 32
features, and the unchanged output head classifies their concatenation
with $h$. This is an architectural use of two observed profiles, not a
learned physical time derivative or evidence of a causal biological mechanism.

The rejected source \texttt{55c3487863c5} adds the existing drug embedding
to the gate input, changing $G$ from a 96-to-32 to a 128-to-32 map and
adding 1,024 weights. The other encoders, latent widths and output head
are unchanged. Its visible AP changes are $-0.021498$ on development and
$-0.024306$ on the gate; the fixed selector therefore retains the previous
model. These are measured diagnostic differences, not significance tests.

The sealed round-3 hypothesis context contains the preceding rejection
and its decision identity. Its new proposal references the decline in both
visible AP measures and proposes dropout of 0.2 in the prediction head.
The next candidate record binds the retained source \texttt{07ba285c6b6c}
to evidence \texttt{578bdbdafc1c}; its own record identity is
\texttt{1f6213b18420}. This verifies that the failed attempt is available
in an actually executed subsequent proposal, while avoiding the stronger
claim that the feedback caused that proposal or improved final outcomes.
